\section{Literature Review}

\subsection{How Chess Is Studied in Data Science and Economics}

Chess has several properties that make it unusually attractive for empirical economics and data science. Players are different in terms of observable skill, they face repeated high-stakes decisions, outcomes are objectively measured, and the full sequence of choices is recorded. In addition, modern chess engines provide a computational benchmark for decision quality. This combination is rare. Many economic settings have high stakes but incomplete information on decision quality. Many laboratory tasks have clean measurement but limited external validity. Chess sits between these extremes: it is a real competitive market for skill, but it also produces structured data.

Early economics work used chess players as highly skilled decision makers. \citet{levitt2011checkmate} use world-class chess players to study backward induction. Their paper is important for this thesis because it treats chess players as expert strategic agents while still testing whether expert status transfers to other decision tasks. The lesson is that chess skill is not simply a general index of rationality. Instead, chess expertise is domain-specific and should be studied with attention to the exact decision environment. This is directly relevant here because a player who is strong in classical or increment chess may not be equally advantaged in a no-increment online format.

Another stream of work uses chess to study risk, patience, and gendered strategic behavior. \citet{gerdes2010strategic} analyze a large panel of expert chess games and find gender differences in strategic risk-taking, including evidence that male players choose more aggressive openings against female opponents. \citet{gransmark2012masters} studies impatience and self-control in high-level chess games, linking game length and behavior around time controls to economic concepts of impatience. These papers motivate the style and risk sections of this thesis. A no-increment time-control shock may not affect only average accuracy; it can change which strategic profiles are rewarded.

\subsection{Chess, Time Pressure, and Cognitive Performance}

The closest conceptual literature studies chess as a high-resolution environment for cognitive performance under constraints. \citet{strittmatter2020lifecycle} use more than 1.6 million moves from professional games and engine evaluations to estimate life-cycle patterns in cognitive performance. Their main result is a hump-shaped age-performance profile. They also control for features such as color, game length, complexity, period, cohort, and player fixed effects. This paper is closely related in measurement strategy because it uses engine-based move quality to study how age-related performance changes after a rule change. The difference is that the present setting focuses on a platform reform in elite online blitz rather than long-run over-the-board chess.

\citet{kuenn2022remote} use professional chess to study cognitive performance in remote work. Their design compares move quality in remote and in-person settings and uses chess engines to evaluate decision quality. This paper is important for two reasons. First, it shows that chess can be used to study broader economic questions about remote work and cognitive production. Second, it demonstrates that move-level engine measures can serve as a credible performance metric when the task itself has a clear computational benchmark. The current thesis follows that logic but focuses on time-control design rather than location of work.

\citet{strittmatter2022speed} study speed, quality, and the timing of complex chess decisions. Their contribution is particularly relevant because their framework shows that decision speed, remaining time, complexity, and move quality are jointly important. This thesis follows the same logic in a different institutional setting: PGN clock annotations are used to recover remaining time and move-level time spent, while Stockfish evaluations measure the quality cost of decisions made under different clock states.

\citet{zegners2020bounded} propose a computational benchmarking approach to bounded rationality in chess. They compare human moves to engine benchmarks and link deviations to position characteristics, fatigue, time pressure, and complexity. This paper is conceptually close to the move-level part of the thesis. The present contribution is not to estimate a general model of bounded rationality, but to use a platform rule change to ask whether the value of different kinds of chess skill changed when the time-control technology changed.

\subsection{Gender, Competition, and Tournament Pressure}

Chess has also been used to study gender and competition. \citet{dilmaghani2020gender} studies gender differences under time constraints using FIDE tournament data. \citet{backus2023gender} analyze gender composition and performance in chess games using a measure of within-game quality of play. They find that women make more mistakes when playing against men, while men do not play worse against women. These papers motivate several hypotheses in the present thesis: whether female players were differentially affected by the 5+0 reform, whether female-male differences appear specifically in late or critical positions, and whether mixed-gender games changed after the reform.

The results in this thesis are more cautious than a simple gender story. The broad female-status interactions on result and accuracy are not robust. A mixed-gender within-game result is statistically strong, but it is not matched by a clear accuracy effect and applies to a selected subset of games. Move-level female interactions do not show a broad female-specific late or critical-position deterioration. This distinction matters because the prior literature shows that chess can reveal gendered behavior, but it does not imply that every time-pressure or rule-change shock will generate robust gender heterogeneity.

Tournament pressure is another relevant theme. Titled Tuesday is not only a sequence of independent games; it is an 11-round Swiss tournament with cash prizes, rank pressure, late-round incentives, and elimination dynamics. The thesis therefore studies bubble players, eliminated players, pairing paths, final rank, completion, and participation. The results show large bubble and eliminated-zone interactions, but placebo and pretrend checks imply that these threshold dynamics were partly present before the reform. The most robust tournament-design result is instead the schedule displacement effect: old late-slot regulars are less likely to participate and less likely to complete the tournament after the second weekly slot disappears.

\subsection{Chess in Data Science and Computer Science}

Chess has a separate and older role in computer science. It is one of the canonical domains in which researchers studied search, planning, evaluation functions, and artificial intelligence. \citet{shannon1950programming} framed computer chess around game-tree search, position evaluation, and selective pruning. \citet{newell1958chess} used chess-playing programs to discuss computational complexity and heuristic problem solving. Later, \citet{campbell2002deepblue} describe Deep Blue, the chess system that defeated Garry Kasparov in 1997 and became a public milestone for search-based artificial intelligence. These classical papers matter for this thesis because they explain why chess became a natural benchmark for algorithmic decision quality: every move is a discrete action in a fully observed state space, and the quality of a move can be approximated by sufficiently strong search and evaluation.

The last decade changed the computer-chess literature in two ways. The first change is the move from hand-engineered evaluation and search toward deep learning and reinforcement learning. \citet{david2016deepchess} train a deep neural system on large chess-position datasets, showing that neural representations can learn useful chess evaluation features from data. \citet{silver2018alphazero} then show that AlphaZero can master chess, shogi, and Go through self-play reinforcement learning, without relying on human games as training examples. \citet{mcgrath2022alphazero} study AlphaZero's internal representations and find that many human chess concepts emerge during training. \citet{maharaj2022chessai} compare competing modern chess-engine paradigms, especially Stockfish-style search and evaluation against neural-network-based engines such as Leela Chess Zero. These papers are important for empirical work because they make engine evaluations credible as a high-quality computational benchmark for human move quality.

The second recent change is that chess is increasingly used as a data-science environment for modeling human behavior, not only for building stronger engines. \citet{toshniwal2022chess} use chess move notation as a controlled test environment for language-model state tracking. Their key idea is that chess is textual, sequential, deterministic, and externally verifiable: a model trained only on move strings can be checked against the true board state and legal move set. This is closely related to the data structure in the present thesis. PGN files are not just text; after parsing, they become sequences of legally valid board states, moves, and engine-evaluated decisions.

Recent human-AI work also uses chess to predict what people actually do, rather than what an optimal engine would do. \citet{mcilroyyoung2020maia} introduce Maia, a neural chess model designed to predict human moves at different skill levels. \citet{mcilroyyoung2022individual} extend this idea to individual players by fine-tuning models to capture player-specific behavior and by using chess moves for behavioral stylometry. This line of work is directly connected to the style component of this paper. Instead of treating every deviation from the engine as the same kind of error, it asks whether move sequences reveal stable and measurable differences across players.

For data science, chess therefore has two useful layers. At the objective layer, engines such as Stockfish transform raw moves into centipawn loss, blunders, critical positions, conversion states, and recovery states. We use in this work Stockfish outputs to better understand what happened in this game, move by move. At the behavioral layer, move sequences can be transformed into player representations, style clusters, and interpretable features. The present thesis uses both layers. Stockfish evaluations provide transparent measures of decision quality, while pre-change move sequences provide style features and neural embeddings. The point is not to build a stronger chess engine, but to use computer-science tools to measure how human decision making changed after the Titled Tuesday time-control reform.

The thesis deliberately combines neural style embeddings with interpretable style indexes. Neural embeddings can capture rich behavioral patterns, but they are hard to interpret in economic terms. Interpretable features make the style analysis more transparent: clean style, chaos creation, self-risk, practical chaos, defensive skill, complexity tolerance, conversion skill, and opening specialization. The style results therefore use the two approaches differently. The interpretable indexes identify relative score penalties for historically chaotic profiles after the switch to 5+0, while the neural embeddings show that stable style clusters can be recovered from pre-change moves and that all recovered clusters improved in raw move accuracy after the reform.

\subsection{Gap in the Existing Literature Filled by This Work}

The existing literature establishes that chess is useful for studying experts decision making, gendered competition, cognitive performance, remote work, time pressure, bounded rationality and even human-AI modeling. This paper contributes a different empirical object: a large online platform rule change in a recurring elite tournament. The setting is valuable because the same players appear repeatedly, the reform is sharp in time, and the change affects both the time-control technology and the participation schedule.Conceptually, this resembles a regression discontinuity design, since the \textit{format\_5\_0} coefficient - which accounts for the change following the transition to the new rules—precisely captures the discontinuity at the point where the rules changed.

The work also contributes methodologically by linking three levels of analysis. The first level is standard panel econometrics/data science using player-game outcomes, ratings, accuracy, demographics, and fixed effects. The second level is move-level engine and clock analysis using Stockfish centipawn loss, blunders, conversion, defensive recovery, critical positions, remaining time, low-clock states, and time spent. The third level is player style measurement using pre-change interpretable features and neural contrastive embeddings. The combination is central: game-level regressions show that the return to skill increased, while move-level and clock results explain where that skill premium appears inside games.

The remaining measurement caveat is more specific. The clock variables are recovered from PGN clock annotations and are therefore measures of recorded remaining time, not direct measures of attention, stress, or intended thinking effort. Very fast online actions, premoves, interface latency, and network conditions may affect the exact mapping from clock change to cognitive decision time. For this reason, the thesis interprets clock coefficients as evidence about the tournament's time-budget technology rather than as pure psychological estimates of pressure.
